\Section{Weakest-Precondition Analysis}
\label{sec:wp}

We run a textbook weakest-precondition (\WP) analysis~\cite{dijkstra1975wp,BL05} over the same bytecode that \MVP normally hands to its Boogie back-end, and read off abort conditions, modifies clauses, and the postconditions that are direct consequences of the function body.
The novelty here is not in the \WP calculus, but in the \emph{shape of the bytecode it operates on}: by the time the analysis runs, \MVP's standard transformation pipeline has already eliminated all Move references and injected every available specification as in-line \texttt{assume}/\texttt{assert} statements.
The result is a reference-free, contract-annotated control flow on which backward predicate propagation reduces to the classical case.

\Paragraph{Reference elimination}
Move has Rust-style references: |&T| for shared reads and |&mut T| for exclusive writes, both with a stack-bounded lifetime.
Predicate-transformer reasoning over such references is awkward, because the meaning of |*r = v| depends on which root location |r| was derived from.
\MVP avoids this by rewriting bytecode into a reference-free intermediate form~\cite{DillGPQXZ22}.
Shared references are inlined: |&T| disappears and reads through it become reads of the underlying value.
Mutable references are replaced by in/out parameters.

\Paragraph{Specification instrumentation}
Before \WP runs, the bytecode will run through the regular process of specification instrumentation.
Each function-level |requires|, |aborts_if|, |ensures|, or |modifies| clause, and each loop |invariant|, is compiled into an |assume| or |assert| at the appropriate program point.
Calls are instrumented uniformly via the \emph{behavioral predicates} introduced for higher-order functions in~\cite{GrieskampEtAl26HO}: |requires_of<f>|, |aborts_of<f>|, |ensures_of<f>|, |modifies_of<f>|, and |result_of<f>| reify a callee's contract as predicates over the call site's arguments and result.
Loops are broken at their headers via the standard \emph{cut-point} encoding of Barnett and Leino for unstructured programs~\cite{BL05}, with the loop |invariant| playing the role of cut-point predicate: the back-edge is severed, the invariant is asserted before and after each iteration, and the loop-modified locations are havoced in between, leaving an acyclic fragment through which \WP propagates as through any straight-line code.
This is the one place where the construction has a genuine \emph{prerequisite}: without an inductive invariant the inferred specification past the loop becomes vacuous, and an inductive invariant cannot in general be recovered from the loop body alone --- this is where AI-proposed candidates enter the picture (Sec.~\ref{sec:skills}).
At the input to \WP, the bytecode therefore looks like a sequence of pure assignments interleaved with |assume| and |assert| over plain values; the surface-level features (references, resources, calls, loops) are gone.
AI-proposed clauses (loop invariants, conservation laws) enter via the same instrumentation pass, so a subsequent \WP run propagates through them as if they had been written by hand.

%%% Local Variables:
%%% mode: latex
%%% TeX-master: "main"
%%% End:
